\begin{frame}{Gap Identified from Related Work}
\begin{center}
\textbf{The reviewed approaches leave limitations when station and flow fairness are considered together.}
\end{center}

\begin{columns}[T,onlytextwidth]
\column{0.32\textwidth}
\textbf{MAC approaches}
\begin{itemize}
\item EIFS, back-off, and CW adaptation.
\item Some need topology-dependent values; most per-flow CW studies are single-hop.
\end{itemize}

\column{0.32\textwidth}
\textbf{TCP approaches}
\begin{itemize}
\item NRED and TCP window or rate adjustment.
\item Primarily control TCP behavior.
\end{itemize}

\column{0.32\textwidth}
\textbf{Cross-layer approaches}
\begin{itemize}
\item Use MAC information to adapt TCP.
\item A cited method omits MAC contention and cannot exactly identify intermediate-node efficiency.
\end{itemize}
\end{columns}

\vspace{-0.05cm}
\begin{center}
\begin{tikzpicture}[node distance=0.55cm]
\node[draw=MethodColor,rounded corners,fill=CardBg,text width=4.2cm,align=center,inner ysep=3pt] (cw) {Suitable CW for each station};
\node[draw=MethodColor,rounded corners,fill=CardBg,text width=4.2cm,align=center,inner ysep=3pt,right=of cw] (rate) {TCP sending rate for each flow};
\node[draw=WarningColor,rounded corners,fill=CardBg,right=of rate] (catra) {CATRA};
\draw[-{Latex[length=2.2mm]},draw=MethodColor] (cw)--(rate);
\draw[-{Latex[length=2.2mm]},draw=MethodColor] (rate)--(catra);
\end{tikzpicture}
\end{center}
\end{frame}
